% problems1.tex — Problem Set and Answers for Chapter 1

\begin{problemset}

\subsection*{Analytical Exercises}

\exercise[1] (Proof that $\vec{b}$ minimizes $SSR$)\quad
Let $\vec{b}$ be the OLS estimator of $\bm{\beta}$. Prove that, for
any hypothetical estimate, $\widetilde{\bm{\beta}}$, of $\bm{\beta}$,
%
\begin{equation}
(\vec{y} - \vec{X}\widetilde{\bm{\beta}})'(\vec{y} - \vec{X}\widetilde{\bm{\beta}})
\geq (\vec{y} - \vec{X}\vec{b})'(\vec{y} - \vec{X}\vec{b}).
\label{eq:ps1.1}
\end{equation}
%
In your proof, use the add-and-subtract strategy: take
$\vec{y} - \vec{X}\widetilde{\bm{\beta}}$, add $\vec{X}\vec{b}$ to it and
then subtract the same from it. It produces the decomposition of
$\vec{y} - \vec{X}\widetilde{\bm{\beta}}$:
%
\[
\vec{y} - \vec{X}\widetilde{\bm{\beta}}
= (\vec{y} - \vec{X}\vec{b}) + (\vec{X}\vec{b} - \vec{X}\widetilde{\bm{\beta}}).
\]
%
\textbf{Hint:}
$(\vec{y} - \vec{X}\widetilde{\bm{\beta}})'(\vec{y} - \vec{X}\widetilde{\bm{\beta}})
= [(\vec{y} - \vec{X}\vec{b}) + \vec{X}(\vec{b} - \widetilde{\bm{\beta}})]'
[(\vec{y} - \vec{X}\vec{b}) + \vec{X}(\vec{b} - \widetilde{\bm{\beta}})]$.
Using the normal equations, show that this equals
%
\[
(\vec{y} - \vec{X}\vec{b})'(\vec{y} - \vec{X}\vec{b})
+ (\vec{b} - \widetilde{\bm{\beta}})'\vec{X}'\vec{X}(\vec{b} - \widetilde{\bm{\beta}}).
\]

\exercise[2] (The annihilator associated with the vector of ones)\quad
Let $\vec{1}$ be the $n$-dimensional column vector of ones, and let
$\vec{M}_1 \equiv \vec{I}_n - \vec{1}(\vec{1}'\vec{1})^{-1}\vec{1}'$.
That is, $\vec{M}_1$ is the annihilator associated with $\vec{1}$.
Prove the following:

\begin{enumerate}[label=(\alph*)]
\item $\vec{M}_1$ is symmetric and idempotent.

\item $\vec{M}_1 \vec{1} = \vec{0}$.

\item $\vec{M}_1 \vec{y} = \vec{y} - \bar{y} \cdot \vec{1}$ where
\[
\bar{y} = \frac{1}{n} \sum_{i=1}^{n} y_i.
\]
$\vec{M}_1 \vec{y}$ is the vector of \textbf{deviations from the mean}.

\item $\vec{M}_1 \vec{X} = \vec{X} - \vec{1}\bar{\vec{x}}'$ where
$\bar{\vec{x}} = \vec{X}'\vec{1}/n$. The $k$-th element of the
$K \times 1$ vector $\bar{\vec{x}}$ is
$\frac{1}{n}\sum_{i=1}^{n} x_{ik}$.
\end{enumerate}

\exercise[3] (Deviation-from-the-mean regression)\quad
Consider a regression model with a constant. Let $\vec{X}$ be partitioned as
%
\[
\underset{(n \times K)}{\vec{X}}
= \begin{bmatrix}
\underset{n \times 1}{\vec{1}} & \vdots & \underset{n \times (K-1)}{\vec{X}_2}
\end{bmatrix}
\]
%
so the first regressor is a constant. Partition $\bm{\beta}$ and $\vec{b}$ accordingly:
%
\[
\bm{\beta} = \begin{bmatrix} \beta_1 \\ \bm{\beta}_2 \end{bmatrix}
\begin{array}{l} \leftarrow \text{scalar} \\ \leftarrow (K-1) \times 1 \end{array}, \quad
\vec{b} = \begin{bmatrix} b_1 \\ \vec{b}_2 \end{bmatrix}.
\]
%
Also let $\widetilde{\vec{X}}_2 \equiv \vec{M}_1 \vec{X}_2$ and
$\tilde{\vec{y}} \equiv \vec{M}_1 \vec{y}$. They are the deviations from the mean for
the nonconstant regressors and the dependent variable. Prove the following:

\begin{enumerate}[label=(\alph*)]
\item The $K$ normal equations are
%
\[
\bar{y} - b_1 - \bar{\vec{x}}_2'\vec{b}_2 = 0
\]
%
where $\bar{\vec{x}}_2 = \vec{X}_2'\vec{1}/n$,
%
\[
\vec{X}_2'\vec{y} - n \cdot b_1 \cdot \bar{\vec{x}}_2
- \vec{X}_2'\vec{X}_2 \vec{b}_2
= \underset{((K-1)\times 1)}{\vec{0}}.
\]

\item $\vec{b}_2 = (\widetilde{\vec{X}}_2'\widetilde{\vec{X}}_2)^{-1}
\widetilde{\vec{X}}_2'\tilde{\vec{y}}$.
\textbf{Hint:} Substitute the first normal equation into the other
$K - 1$ equations to eliminate $b_1$ and solve for $\vec{b}_2$.
This is a generalization of the result you proved in Review Question~3
in Section~1.2.
\end{enumerate}

\exercise[4] (Partitioned regression, generalization of Exercise~3)\quad
Let $\vec{X}$ be partitioned as
%
\[
\underset{(n \times K)}{\vec{X}}
= \begin{bmatrix}
\underset{(n \times K_1)}{\vec{X}_1} & \vdots &
\underset{(n \times K_2)}{\vec{X}_2}
\end{bmatrix}.
\]
%
Partition $\bm{\beta}$ accordingly:
%
\[
\bm{\beta} = \begin{bmatrix} \bm{\beta}_1 \\ \bm{\beta}_2 \end{bmatrix}
\begin{array}{l} \leftarrow K_1 \times 1 \\ \leftarrow K_2 \times 1 \end{array}.
\]
%
Thus, the regression can be written as
%
\[
\vec{y} = \vec{X}_1 \bm{\beta}_1 + \vec{X}_2 \bm{\beta}_2 + \bm{\varepsilon}.
\]
%
Let $\vec{P}_1 \equiv \vec{X}_1(\vec{X}_1'\vec{X}_1)^{-1}\vec{X}_1'$,
$\vec{M}_1 \equiv \vec{I} - \vec{P}_1$,
$\widetilde{\vec{X}}_2 \equiv \vec{M}_1 \vec{X}_2$ and
$\tilde{\vec{y}} \equiv \vec{M}_1 \vec{y}$.
Thus, $\tilde{\vec{y}}$ is the residual vector from the regression of $\vec{y}$
on $\vec{X}_1$, and the $k$-th column of $\widetilde{\vec{X}}_2$ is the residual
vector from the regression of the corresponding $k$-th column of
$\vec{X}_2$ on $\vec{X}_1$. Prove the following:

\begin{enumerate}[label=(\alph*)]
\item The normal equations are
%
\begin{align}
\vec{X}_1'\vec{X}_1 \vec{b}_1 + \vec{X}_1'\vec{X}_2 \vec{b}_2
&= \vec{X}_1'\vec{y}, \tag{$*$} \label{eq:ps1.4star}\\
\vec{X}_2'\vec{X}_1 \vec{b}_1 + \vec{X}_2'\vec{X}_2 \vec{b}_2
&= \vec{X}_2'\vec{y}. \tag{$**$} \label{eq:ps1.4starstar}
\end{align}

\item $\vec{b}_2 = (\widetilde{\vec{X}}_2'\widetilde{\vec{X}}_2)^{-1}
\widetilde{\vec{X}}_2'\tilde{\vec{y}}$.
That is, $\vec{b}_2$ can be obtained by regressing the residuals
$\tilde{\vec{y}}$ on the matrix of residuals $\widetilde{\vec{X}}_2$.
\textbf{Hint:} Derive $\vec{X}_1\bm{\beta}_1 = -\vec{P}_1\vec{X}_2\bm{\beta}_2
+ \vec{P}_1\vec{y}$ from~\eqref{eq:ps1.4star}. Substitute this
into~\eqref{eq:ps1.4starstar} to obtain
$\vec{X}_2'\vec{M}_1\vec{X}_2\bm{\beta}_2 = \vec{X}_2'\vec{M}_1\vec{y}$.
Then use the fact that $\vec{M}_1$ is symmetric and idempotent.
Or, if you wish, you can apply the brute force of the partitioned inverse
formula~(A.10) of Appendix~A to the coefficient matrix
%
\[
\vec{X}'\vec{X} = \begin{bmatrix}
\vec{X}_1'\vec{X}_1 & \vec{X}_1'\vec{X}_2 \\
\vec{X}_2'\vec{X}_1 & \vec{X}_2'\vec{X}_2
\end{bmatrix}.
\]
%
Show that the second diagonal block of $(\vec{X}'\vec{X})^{-1}$ is
$(\widetilde{\vec{X}}_2'\widetilde{\vec{X}}_2)^{-1}$.

\item The residuals from the regression of $\tilde{\vec{y}}$ on
$\widetilde{\vec{X}}_2$ numerically equals $\vec{e}$, the residuals from the
regression of $\vec{y}$ on $\vec{X}$ ($\equiv (\vec{X}_1 \vdots \vec{X}_2)$).
\textbf{Hint:} If $\vec{e}$ is the residual from the regression of $\vec{y}$
on $\vec{X}$,
%
\[
\vec{y} = \vec{X}_1 \vec{b}_1 + \vec{X}_2 \vec{b}_2 + \vec{e}.
\]
%
Premultiplying both sides by $\vec{M}_1$ and using
$\vec{M}_1\vec{X}_1 = \vec{0}$, we obtain
%
\[
\tilde{\vec{y}} = \widetilde{\vec{X}}_2 \vec{b}_2 + \vec{M}_1\vec{e}.
\]
%
Show that $\vec{M}_1\vec{e} = \vec{e}$ and observe that $\vec{b}_2$ equals the
OLS coefficient estimate in the regression of $\tilde{\vec{y}}$ on
$\widetilde{\vec{X}}_2$.

\item $\vec{b}_2 = (\widetilde{\vec{X}}_2'\widetilde{\vec{X}}_2)^{-1}
\widetilde{\vec{X}}_2'\vec{y}$.
Note the difference from~(b). Here, the vector of dependent variable is
$\vec{y}$, not $\tilde{\vec{y}}$. Are the residuals from the regression of
$\vec{y}$ on $\widetilde{\vec{X}}_2$ numerically the same as $\vec{e}$?
[Answer: No.]\enspace
Is the $SSR$ from the regression of $\vec{y}$ on $\widetilde{\vec{X}}_2$ the
same as the $SSR$ from the regression of $\tilde{\vec{y}}$ on
$\widetilde{\vec{X}}_2$? [Answer: No.]
\end{enumerate}

\noindent
The results in (b)--(d) are known as the \textbf{Frisch--Waugh Theorem}.

\begin{enumerate}[label=(\alph*),resume]
\item Show:
%
\[
\tilde{\vec{y}}'\tilde{\vec{y}} - \vec{e}'\vec{e}
= \tilde{\vec{y}}'\vec{X}_2(\vec{X}_2'\vec{M}_1\vec{X}_2)^{-1}\vec{X}_2'\tilde{\vec{y}}.
\]
%
\textbf{Hint:} Apply the general decomposition formula~(1.2.15) to the
regression in~(c) to derive
%
\[
\tilde{\vec{y}}'\tilde{\vec{y}}
= \vec{b}_2'\widetilde{\vec{X}}_2'\widetilde{\vec{X}}_2\vec{b}_2 + \vec{e}'\vec{e}.
\]
%
Then use~(b).

\item Consider the following four regressions:
\begin{enumerate}[label=(\arabic*)]
\item regress $\tilde{\vec{y}}$ on $\vec{X}_1$.
\item regress $\tilde{\vec{y}}$ on $\widetilde{\vec{X}}_2$.
\item regress $\tilde{\vec{y}}$ on $\vec{X}_1$ and $\vec{X}_2$.
\item regress $\tilde{\vec{y}}$ on $\vec{X}_2$.
\end{enumerate}

Let $SSR_j$ be the sum of squared residuals from regression $j$. Show:

\begin{enumerate}[label=(\roman*)]
\item $SSR_1 = \tilde{\vec{y}}'\tilde{\vec{y}}$.
\textbf{Hint:} $\tilde{\vec{y}}$ is constructed so that
$\vec{X}_1'\tilde{\vec{y}} = \vec{0}$, so $\vec{X}_1$ should have
no explanatory power.

\item $SSR_2 = \vec{e}'\vec{e}$. \textbf{Hint:} Use~(c).

\item $SSR_3 = \vec{e}'\vec{e}$. \textbf{Hint:} Apply the Frisch--Waugh
Theorem on regression~(3). $\vec{M}_1\tilde{\vec{y}} = \tilde{\vec{y}}$.

\item Verify by numerical example that $SSR_4$ is not necessarily equal to
$\vec{e}'\vec{e}$.
\end{enumerate}
\end{enumerate}

\exercise[5] (Restricted regression and $F$)\quad
In the restricted least squares, the sum of squared residuals is minimized
subject to the constraint implied by the null hypothesis
$\vec{R}\bm{\beta} = \vec{r}$. Form the Lagrangian as
%
\[
\mathcal{L} = \frac{1}{2}(\vec{y} - \vec{X}\widetilde{\bm{\beta}})'
(\vec{y} - \vec{X}\widetilde{\bm{\beta}})
+ \bm{\lambda}'(\vec{R}\widetilde{\bm{\beta}} - \vec{r}),
\]
%
where $\bm{\lambda}$ here is the $\#\vec{r}$-dimensional vector of Lagrange
multipliers (recall: $\vec{R}$ is $\#\vec{r} \times K$,
$\widetilde{\bm{\beta}}$ is $K \times 1$, and $\vec{r}$ is
$\#\vec{r} \times 1$). Let $\widehat{\bm{\beta}}$ be the restricted least
squares estimator of $\bm{\beta}$. It is the solution to the constrained
minimization problem.

\begin{enumerate}[label=(\alph*)]
\item Let $\vec{b}$ be the unrestricted OLS estimator. Show:
%
\begin{align}
\widehat{\bm{\beta}} &= \vec{b} - (\vec{X}'\vec{X})^{-1}\vec{R}'
[\vec{R}(\vec{X}'\vec{X})^{-1}\vec{R}']^{-1}(\vec{R}\vec{b} - \vec{r}),
\label{eq:ps1.5a1} \\
\bm{\lambda} &= [\vec{R}(\vec{X}'\vec{X})^{-1}\vec{R}']^{-1}
(\vec{R}\vec{b} - \vec{r}). \label{eq:ps1.5a2}
\end{align}
%
\textbf{Hint:} The first-order conditions are
$\vec{X}'\vec{y} - (\vec{X}'\vec{X})\widehat{\bm{\beta}} = \vec{R}'\bm{\lambda}$
or $\vec{X}'(\vec{y} - \vec{X}\widehat{\bm{\beta}}) = \vec{R}'\bm{\lambda}$.
Combine this with the constraint
$\vec{R}\widehat{\bm{\beta}} = \vec{r}$ to solve for $\bm{\lambda}$ and
$\widehat{\bm{\beta}}$.

\item Let $\hat{\bm{\varepsilon}} \equiv \vec{y} - \vec{X}\widehat{\bm{\beta}}$,
the residuals from the restricted regression. Show:
%
\begin{align*}
SSR_R - SSR_U
&= (\vec{b} - \widehat{\bm{\beta}})' (\vec{X}'\vec{X})
   (\vec{b} - \widehat{\bm{\beta}}) \\
&= (\vec{R}\vec{b} - \vec{r})'
   [\vec{R}(\vec{X}'\vec{X})^{-1}\vec{R}']^{-1}
   (\vec{R}\vec{b} - \vec{r}) \\
&= \bm{\lambda}'\vec{R}(\vec{X}'\vec{X})^{-1}\vec{R}'\bm{\lambda} \\
&= \hat{\bm{\varepsilon}}'\vec{P}\hat{\bm{\varepsilon}},
\end{align*}
%
where $\vec{P}$ is the projection matrix. \textbf{Hint:} For the first
equality, use the add-and-subtract strategy:
%
\begin{align*}
SSR_R &= (\vec{y} - \vec{X}\widehat{\bm{\beta}})'
         (\vec{y} - \vec{X}\widehat{\bm{\beta}}) \\
      &= [(\vec{y} - \vec{X}\vec{b})
         + \vec{X}(\vec{b} - \widehat{\bm{\beta}})]'
         [(\vec{y} - \vec{X}\vec{b})
         + \vec{X}(\vec{b} - \widehat{\bm{\beta}})].
\end{align*}
%
Use the normal equations $\vec{X}'(\vec{y} - \vec{X}\vec{b}) = \vec{0}$.
For the second and third equalities, use~(a). To prove the fourth equality,
the easiest way is to use the first-order condition mentioned in~(a) that
$\vec{R}'\bm{\lambda} = \vec{X}'\hat{\bm{\varepsilon}}$.

\item Verify that you have proved in~(b) that $(1.4.9) = (1.4.11)$.
\end{enumerate}

\exercise[6] (Proof of the decomposition~(1.2.17))\quad
Take the unrestricted model to be a regression where one of the regressors
is a constant, and the restricted model to be a regression where the only
regressor is a constant.

\begin{enumerate}[label=(\alph*)]
\item Show that~(b) in the previous exercise is the
decomposition~(1.2.17) for this case. \textbf{Hint:} What is
$\widehat{\bm{\beta}}$ for this case? Show that
$SSR_R = \sum_i (y_i - \bar{y})^2$ and
$(\vec{b} - \widehat{\bm{\beta}})' (\vec{X}'\vec{X})
(\vec{b} - \widehat{\bm{\beta}}) = \sum_i (\hat{y}_i - \bar{y})^2$.

\item ($R^2$ as an $F$-ratio)\quad
For a regression where one of the regressors is a constant, prove that
%
\[
F = \frac{R^2/(K-1)}{(1 - R^2)/(n - K)}.
\]
\end{enumerate}

\exercise[7] (Hausman principle in finite samples)\quad
For the generalized regression model, prove the following. Here, it is
understood that the expectations, variances, and covariances are all
conditional on $\vec{X}$.

\begin{enumerate}[label=(\alph*)]
\item $\Cov(\widehat{\bm{\beta}}_{\text{GLS}},\,
\vec{b} - \widehat{\bm{\beta}}_{\text{GLS}}) = \vec{0}$.
\textbf{Hint:} Recall that, for any two random vectors $\vec{x}$ and
$\vec{y}$,
%
\[
\Cov(\vec{x}, \vec{y})
\equiv \E\bigl[(\vec{x} - \E(\vec{x}))(\vec{y} - \E(\vec{y}))'\bigr].
\]
%
So
%
\[
\Cov(\vec{A}\vec{x}, \vec{B}\vec{y}) = \vec{A}\,\Cov(\vec{x},\vec{y})\,\vec{B}'.
\]
%
Also, since $\bm{\beta}$ is nonrandom,
%
\[
\Cov(\widehat{\bm{\beta}}_{\text{GLS}},\,
\vec{b} - \widehat{\bm{\beta}}_{\text{GLS}})
= \Cov(\widehat{\bm{\beta}}_{\text{GLS}} - \bm{\beta},\,
\vec{b} - \widehat{\bm{\beta}}_{\text{GLS}}).
\]

\item Let $\widetilde{\bm{\beta}}$ be any unbiased estimator and define
$\vec{q} \equiv \widetilde{\bm{\beta}} - \widehat{\bm{\beta}}_{\text{GLS}}$.
Assume $\widetilde{\bm{\beta}}$ is such that
$\vec{V}_{\vec{q}} \equiv \Var(\vec{q})$ is nonsingular. Prove:
$\Cov(\widehat{\bm{\beta}}_{\text{GLS}}, \vec{q}) = \vec{0}$.
(If we set $\widetilde{\bm{\beta}} = \vec{b}$, we are back to~(a).)
\textbf{Hint:} Define:
$\widehat{\bm{\beta}} \equiv \widehat{\bm{\beta}}_{\text{GLS}}
+ \vec{H}\vec{q}$ for some $\vec{H}$. Show:
%
\[
\Var(\widehat{\bm{\beta}})
= \Var(\widehat{\bm{\beta}}_{\text{GLS}})
+ \vec{C}\vec{H}' + \vec{H}\vec{C}'
+ \vec{H}\vec{V}_{\vec{q}}\vec{H}',
\]
%
where $\vec{C} \equiv \Cov(\widehat{\bm{\beta}}_{\text{GLS}}, \vec{q})$.
Show that, if $\vec{C} \neq \vec{0}$ then $\Var(\widehat{\bm{\beta}})$ can
be made smaller than $\Var(\widehat{\bm{\beta}}_{\text{GLS}})$ by setting
$\vec{H} = -\vec{C}\vec{V}_{\vec{q}}^{-1}$. Argue that this is in
contradiction to Proposition~1.7(c).

\item (Optional, only for those who are proficient in linear algebra)\quad
Prove: if the $K$ columns of $\vec{X}$ are characteristic vectors of
$\vec{V}$, then $\vec{b} = \widehat{\bm{\beta}}_{\text{GLS}}$, where
$\vec{V}$ is the $n \times n$ variance-covariance matrix of the
$n$-dimensional error vector $\bm{\varepsilon}$.
(So not all unbiased estimators satisfy the requirement in~(b) that
$\Var(\widetilde{\bm{\beta}} - \widehat{\bm{\beta}}_{\text{GLS}})$ be
nonsingular.)
\textbf{Hint:} For any $n \times n$ symmetric matrix $\vec{V}$, there
exists an $n \times n$ matrix $\vec{H}$ such that
$\vec{H}'\vec{H} = \vec{I}_n$ (so $\vec{H}$ is an orthogonal matrix) and
$\vec{H}'\vec{V}\vec{H} = \bm{\Lambda}$, where $\bm{\Lambda}$ is a diagonal
matrix with the characteristic roots (which are real since $\vec{V}$ is
symmetric) of $\vec{V}$ in the diagonal. The columns of $\vec{H}$ are called
the characteristic vectors of $\vec{V}$. Show that
%
\[
\vec{H}^{-1} = \vec{H}', \quad
\vec{H}'\vec{V}^{-1}\vec{H} = \bm{\Lambda}^{-1}, \quad
\vec{H}'\vec{V}^{-1} = \bm{\Lambda}^{-1}\vec{H}'.
\]
%
Without loss of generality, $\vec{X}$ can be taken to be the first $K$
columns of $\vec{H}$. So $\vec{X} = \vec{H}\vec{F}$, where
%
\[
\underset{(n \times K)}{\vec{F}}
= \begin{bmatrix} \vec{I}_K \\ \vec{0} \end{bmatrix}.
\]
\end{enumerate}

\subsection*{Empirical Exercises}

Read Marc Nerlove, ``Returns to Scale in Electricity Supply'' (except
paragraphs of equations~(6)--(9), the part of section~2 from p.~184 on,
and Appendix~A and~C) before doing this exercise. For 145 electric utility
companies in 1955, the file \texttt{NERLOVE.ASC} has data on the following:

\medskip
\noindent
\begin{tabular}{@{}ll}
Column 1: & total costs (call it $TC$) in millions of dollars \\
Column 2: & output ($Q$) in billions of kilowatt hours \\
Column 3: & price of labor ($PL$) \\
Column 4: & price of fuels ($PF$) \\
Column 5: & price of capital ($PK$).
\end{tabular}
\medskip

\noindent
They are from the data appendix of his article. There are 145 observations,
and the observations are ordered in size, observation~1 being the smallest
company and observation~145 the largest. Using the data transformation
facilities of your computer software, generate for each of the 145 firms
the variables required for estimation. To estimate~(1.7.4), for example,
you need to generate $\log(TC)$, a constant, $\log(Q)$, $\log(PL)$,
$\log(PK)$, and $\log(PF)$, for each of the 145 firms.

\begin{enumerate}[label=(\alph*)]
\item (Data question)\quad Does Nerlove's construction of the price of
capital conform to the definition of the user cost of capital?
\textbf{Hint:} Read Nerlove's Appendix~B.4.

\item Estimate the unrestricted model~(1.7.4) by OLS\@. Can you replicate
the estimates in the text?

\item (Restricted least squares)\quad Estimate the restricted
model~(1.7.6) by OLS\@. To do this, you need to generate a new set of
variables for each of the 145 firms. For example, the dependent variable
is $\log(TC/PF)$, not $\log(TC)$. Can you replicate the estimates in the
text? Can you replicate Nerlove's results? Nerlove's estimate of
$\beta_2$, for example, is 0.721 with a standard error of 0.0174 (the
standard error in his paper is 0.175, but it is probably a typographical
error). Where in Nerlove's paper can you find this estimate? What about the
other coefficients? (Warning: You will not be able to replicate Nerlove's
results precisely. One reason is that he used common rather than natural
logarithms; however, this should affect only the estimated intercept term.
The other reason: the data set used for his results is a corrected version
of the data set published with his article.)

As mentioned in the text, the plot of residuals suggests a nonlinear
relationship between $\log(TC)$ and $\log(Q)$. Nerlove hypothesized that
estimated returns to scale varied with the level of output. Following
Nerlove, divide the sample of 145 firms into five subsamples or groups,
each having 29 firms. (Recall that since the data are ordered by level of
output, the first 29 observations will have the smallest output levels,
whereas the last 29 observations will have the largest output levels.)
Consider the following three generalizations of the model~(1.7.6):

\medskip
\noindent
\textit{Model~1:}\enspace Both the coefficients ($\beta$'s) and the error
variance in~(1.7.6) differ across groups.

\noindent
\textit{Model~2:}\enspace The coefficients are different, but the error
variance is the same across groups.

\noindent
\textit{Model~3:}\enspace While each group has common coefficients for
$\beta_3$ and $\beta_4$ (price elasticities) and common error variance, it
has a different intercept term and a different $\beta_2$. Model~3 is what
Nerlove called the hypothesis of neutral variations in returns to scale.

\medskip
For Model~1, the coefficients and error variances specific to groups can be
estimated from
%
\[
\vec{y}^{(j)} = \vec{X}^{(j)} \bm{\beta}^{(j)} + \bm{\varepsilon}^{(j)}
\quad (j = 1, \ldots, 5),
\]
%
where $\vec{y}^{(j)}$ ($29 \times 1$) is the vector of the values of the
dependent variable for group~$j$, $\vec{X}^{(j)}$ ($29 \times 4$) is the
matrix of the values of the four regressors for group~$j$,
$\bm{\beta}^{(j)}$ ($4 \times 1$) is the coefficient vector for group~$j$,
and $\bm{\varepsilon}^{(j)}$ ($29 \times 1$) is the error vector. The second
column of $\vec{X}^{(5)}$, for example, is $\log(Q)$ for
$i = 117, \ldots, 145$. Model~1 assumes conditional homoskedasticity
$\E(\bm{\varepsilon}^{(j)}\bm{\varepsilon}^{(j)\prime} \mid \vec{X}^{(j)})
= \sigma_j^2 \vec{I}_{29}$ within (but not necessarily across) groups.

\item Estimate Model~1 by OLS\@. How well can you replicate Nerlove's
reported results? On the basis of your estimates of $\beta_2$, compute the
point estimates of returns to scale in each of the five groups. What is the
general pattern of estimated scale economies as the level of output
increases? What is the general pattern of the estimated error variance as
output increases?

\medskip
Model~2 assumes for Model~1 that $\sigma_j^2 = \sigma^2$ for all $j$.
This equivariance restriction can be incorporated by stacking vectors and
matrices as follows:
%
\[
\vec{y} = \vec{X}\bm{\beta} + \bm{\varepsilon},
\]
%
where
%
\begin{equation}
\underset{(145 \times 1)}{\vec{y}}
= \begin{bmatrix} \vec{y}^{(1)} \\ \vdots \\ \vec{y}^{(5)} \end{bmatrix}, \quad
\underset{(145 \times 20)}{\vec{X}}
= \begin{bmatrix} \vec{X}^{(1)} & & \\ & \ddots & \\ & & \vec{X}^{(5)} \end{bmatrix}, \quad
\underset{(145 \times 1)}{\bm{\varepsilon}}
= \begin{bmatrix} \bm{\varepsilon}^{(1)} \\ \vdots \\ \bm{\varepsilon}^{(5)} \end{bmatrix}.
\tag{$*$} \label{eq:ps1.emp.star}
\end{equation}
%
In particular, $\vec{X}$ is now a block-diagonal matrix. The equivariance
restriction can be expressed as
$\E(\bm{\varepsilon}\bm{\varepsilon}' \mid \vec{X}) = \sigma^2 \vec{I}_{145}$.
There are now 20 variables derived from the original four regressors. The
145 dimensional vector corresponding to the second variable, for example,
has $\log(Q_1), \ldots, \log(Q_{29})$ as the first 29 elements and zeros
elsewhere. The vector corresponding to the 6th variable, which represents
log output for the second group of firms, has $\log(Q_{30}), \ldots,
\log(Q_{58})$ for the 30th through 58th elements and zeros elsewhere, and
so on.

The stacking operation needed to form the $\vec{y}$ and $\vec{X}$
in~\eqref{eq:ps1.emp.star} can be done easily if your computer software is
matrix-based. Otherwise, you trick your software into accomplishing the same
thing by the use of dummy variables. Define the $j$-th dummy variable as
%
\[
D_{ji} = \begin{cases}
1 & \text{if firm } i \text{ belongs to the } j\text{-th group}, \\
0 & \text{otherwise},
\end{cases}
\quad (i = 1, \ldots, 145).
\]
%
Then the second regressor is $D_{1i} \cdot \log(Q_i)$. The 6th variable is
$D_{2i} \cdot \log(Q_i)$, and so forth.

\item Estimate Model~2 by OLS\@. Verify that the OLS coefficient estimates
here are the same as those in~(d). Also verify that
%
\[
\sum_{j=1}^{5} SSR_j = SSR,
\]
%
where $SSR_j$ is the $SSR$ from the $j$-th group in your estimation of
Model~1 in~(d) and $SSR$ is the $SSR$ from Model~2. This agreement is not
by accident, i.e., not specific to the present data set. Prove that this
agreement for the coefficients and the $SSR$ holds in general, temporarily
assuming just two groups without loss of generality. \textbf{Hint:} First
show that the coefficient estimate is the same between Model~1 and Model~2.
Use formulas~(A.4), (A.5), and~(A.9) of Appendix~A.

\item (Chow test)\quad Model~2 is more general than Model~(1.7.6) because
the coefficients can differ across groups. Test the null hypothesis that the
coefficients are the same across groups. How many equations (restrictions)
are in the null hypothesis? This test is sometimes called the \textbf{Chow
test for structural change}. Calculate the $p$-value of the $F$-ratio.
\textbf{Hint:} This is a linear hypothesis about the coefficients of
Model~2. So take Model~2 to be the maintained hypothesis and~(1.7.6) to be
the restricted model. Use the formula~(1.4.11) for the $F$-ratio.

\medskip
\noindent
\textit{Gauss Tip:}\enspace If $x$ is the $F$-ratio, the Gauss command
\texttt{cdffc(x,df1,df2)} gives the area to the right of $F$ for the $F$
distribution with $df1$ and $df2$ degrees of freedom.

\medskip
\noindent
\textit{TSP Tip:}\enspace The TSP command to do the same is
\texttt{cdf(f, df1=\textit{df1}, df2=\textit{df2}) x}. An output of TSP's
OLS command, \texttt{OLSQ}, is \texttt{@SSR}, which is the $SSR$ for the
regression.

\medskip
\noindent
\textit{RATS Tip:}\enspace The RATS command is
\texttt{cdf ftest x df1 df2}. An output of RATS's OLS command,
\texttt{LINREG}, is \texttt{\%RSS}, which is the $SSR$ for the regression.

\medskip
The restriction in Model~3 that the price elasticities are the same across
firm groups can be imposed on Model~2 by applying the dummy variable
transformation only to the constant and log output. Thus, there are
$12$ ($= 2 \times 5 + 2$) variables in $\vec{X}$. Now $\vec{X}$ looks like
%
\begin{equation}
\vec{X} =
\begin{bmatrix}
1      & \log(Q_1)      & 0 & 0 & \log(PL_1/PF_1)       & \log(PK_1/PF_1) \\
\vdots & \vdots         & \vdots & \vdots & \vdots        & \vdots \\
1      & \log(Q_{29})   & 0 & 0 & \log(PL_{29}/PF_{29})  & \log(PK_{29}/PF_{29}) \\
       &                & \ddots &        & \vdots        & \vdots \\
0      & 0              & 1 & \log(Q_{117})  & \log(PL_{117}/PF_{117}) & \log(PK_{117}/PF_{117}) \\
\vdots & \vdots         & \vdots & \vdots & \vdots        & \vdots \\
0      & 0              & 1 & \log(Q_{145})  & \log(PL_{145}/PF_{145}) & \log(PK_{145}/PF_{145})
\end{bmatrix}
\tag{$**$} \label{eq:ps1.emp.starstar}
\end{equation}

\item Estimate Model~3. The model is a special case of Model~2, with the
hypothesis that the two price elasticities are the same across the five
groups. Test the hypothesis at a significance level of 5 percent, assuming
normality. (Note: Nerlove's $F$-ratio on p.~183 is wrong.)

\medskip
As has become clear from the plot of residuals in Figure~1.7, the
conditional second moment $\E(\varepsilon_i^2 \mid \vec{X})$ is likely to
depend on log output, which is a violation of the conditional
homoskedasticity assumption. This time we do not attempt to test conditional
homoskedasticity, because to do so requires large sample theory and is
postponed until the next chapter. Instead, we pretend to know the form of the
function linking the conditional second moment to log output. The function,
specified below, implies that the conditional second moment varies
continuously with output, contrary to the three models we have considered
above. Also contrary to those models, we assume that the degree of returns
to scale varies continuously with output by including the square of log
output.\footnote{We have derived the log-linear cost function from the
Cobb--Douglas production function. Does there exist a production function
from which this generalized cost function with a quadratic term in log output
can be derived? This is a question of the ``integrability'' of cost functions
and is discussed in detail in Christensen et al.\ (1973).}
Model~4 is

\medskip
\noindent
\textit{Model~4:}
%
\begin{align*}
\log\!\left(\frac{TC_i}{p_{i3}}\right)
&= \beta_1 + \beta_2\,\log(Q_i) + \beta_3\,[\log(Q_i)]^2 \\
&\quad + \beta_4\,\log\!\left(\frac{p_{i1}}{p_{i3}}\right)
       + \beta_5\,\log\!\left(\frac{p_{i2}}{p_{i3}}\right)
       + \varepsilon_i \\[6pt]
\E(\varepsilon_i^2 \mid \vec{X})
&= \sigma^2 \cdot \left(0.0565 + \frac{2.1377}{Q_i}\right)
\quad (i = 1, 2, \ldots, 145)
\end{align*}
%
for some unknown $\sigma^2$.

\item Estimate Model~4 by weighted least squares on the whole sample of
145 firms. (Be careful about the treatment of the intercept; in the equation
after weighting, none of the regressors is a constant.) Plot the residuals.
Is there still evidence for conditional homoskedasticity or further
nonlinearities?
\end{enumerate}

\subsection*{Monte Carlo Exercises}

\textbf{Monte Carlo} analysis simulates a large number of samples from the
model to study the finite-sample distribution of estimators. In this
exercise, we use the technique to confirm the two finite-sample results of
the text: the unbiasedness of the OLS coefficient estimator and the
distribution of the $t$-ratio. The model is the following simple regression
model satisfying Assumptions~1.1--1.5 with $n = 32$. The regression equation
is
%
\begin{align*}
y_i &= \beta_1 + \beta_2 x_i + \varepsilon_i
\quad (i = 1, 2, \ldots, n) \\
\text{or}\quad
\vec{y} &= \vec{1} \cdot \beta_1 + \vec{x} \cdot \beta_2 + \bm{\varepsilon}
= \vec{X}\bm{\beta} + \bm{\varepsilon}, \tag{$*$}
\end{align*}
%
where $\vec{X} = (\vec{1} \vdots \vec{x})$ and
$\bm{\beta} = (\beta_1, \beta_2)'$. The model parameters are
$(\beta_1, \beta_2, \sigma^2)$. As mentioned in the text, a model is a set
of joint distributions of $(\vec{y}, \vec{X})$. We pick a particular joint
distribution by specifying the regression model as follows. Set
$\beta_1 = 1$, $\beta_2 = 0.5$, and $\sigma^2 = 1$. The distribution of
$\vec{x} = (x_1, x_2, \ldots, x_n)'$ is specified by the following AR(1)
process:
%
\begin{equation}
x_i = c + \phi\, x_{i-1} + \eta_i \quad (i = 1, 2, \ldots, n),
\tag{$**$}
\end{equation}
%
where $\{\eta_i\}$ is i.i.d.\ $N(0, 1)$ and
%
\[
x_0 \sim N\!\left(\frac{c}{1 - \phi},\; \frac{1}{1 - \phi^2}\right),
\quad c = 2, \quad \phi = 0.6.
\]
%
This fixes the joint distribution of $(\vec{y}, \vec{X})$. From this
distribution, a large number of samples will be drawn.

In programming the simulation, the following expression for $\vec{x}$ will
be useful. Solve the first-order difference equation~($**$) to obtain
%
\begin{align*}
x_i &= \phi^i x_0 + (1 + \phi + \phi^2 + \cdots + \phi^{i-1})c \\
    &\quad + (\eta_i + \phi\,\eta_{i-1} + \phi^2\,\eta_{i-2}
       + \cdots + \phi^{i-1}\eta_1),
\end{align*}
%
or, in matrix notation,
%
\begin{equation}
\underset{(n \times 1)}{\vec{x}}
= \underset{(n \times 1)}{\vec{r}} \cdot x_0
+ \underset{(n \times 1)}{\vec{d}}
+ \underset{(n \times n)}{\vec{A}}\;
  \underset{(n \times 1)}{\bm{\eta}},
\tag{$***$} \label{eq:ps1.mc.stars}
\end{equation}
%
where $\vec{d} = (d_1, d_2, \ldots, d_n)'$ and
%
\[
d_1 = c, \quad d_2 = (1 + \phi)c, \quad \ldots, \quad
d_i = (1 + \phi + \phi^2 + \cdots + \phi^{i-1})c, \quad \ldots,
\]
%
\[
\vec{r} = \begin{bmatrix} \phi \\ \phi^2 \\ \vdots \\ \phi^n \end{bmatrix},
\quad
\vec{A} = \begin{bmatrix}
1         & 0            & \cdots  & \cdots & 0 \\
\phi      & 1            & 0       & \cdots & 0 \\
\phi^2    & \phi         & 1       & \ddots & 0 \\
\vdots    & \vdots       & \ddots  & \ddots & \vdots \\
\phi^{n-1}& \phi^{n-2}   & \cdots  & \phi   & 1
\end{bmatrix},
\quad
\bm{\eta} = \begin{bmatrix} \eta_1 \\ \eta_2 \\ \vdots \\ \eta_n \end{bmatrix}.
\]
%
\textit{Gauss Tip:}\enspace To form the $\vec{r}$ matrix, use \texttt{seqm}.
To form the $\vec{A}$ matrix, use \texttt{toeplitz} and \texttt{lowmat}.

\begin{enumerate}[label=(\alph*)]
\item Run two Monte Carlo simulations. The first simulation calculates
$\E(\vec{b} \mid \vec{x})$ and the distribution of the $t$-ratio as a
distribution conditional on $\vec{x}$. A computer program for the first
simulation should consist of the following steps.

\begin{enumerate}[label=(\arabic*)]
\item (Generate $\vec{x}$ just once)\quad Using the random number generator,
draw a vector $\bm{\eta}$ of $n$ i.i.d.\ random variables from $N(0, 1)$ and
$x_0$ from $N(c/(1 - \phi),\; 1/(1 - \phi^2))$, and calculate $\vec{x}$
by~\eqref{eq:ps1.mc.stars}. (Calculation of $\vec{x}$ can also be
accomplished recursively by~($**$) with a do loop, but vector operations such
as~\eqref{eq:ps1.mc.stars} consume less CPU time than do loops. This becomes
a consideration in the second simulation, where $\vec{x}$ has to be generated
in each replication.)

\item Set a counter to zero. The counter will record the incidence that
$|t| > t_{0.025}(n - 2)$. Also, set a two-dimensional vector at zero; this
vector will be used for calculating the mean of the OLS estimator $\vec{b}$
of $(\beta_1, \beta_2)'$.

\item Start a do loop of a large number of replications (1 million, say).
In each replication, do the following.

\begin{enumerate}[label=(\roman*)]
\item (Generate $\vec{y}$)\quad Draw an $n$ dimensional vector
$\bm{\varepsilon}$ of $n$ i.i.d.\ random variables from $N(0, 1)$, and
calculate $\vec{y} = (y_1, \ldots, y_n)'$ by~($*$). This $\vec{y}$ is
paired with the same $\vec{x}$ from step~(1) to form a sample
$(\vec{y}, \vec{x})$.

\item From the sample, calculate the OLS estimator $\vec{b}$ and the
$t$-value for $\mathrm{H}_0\colon \beta_2 = 0.5$.

\item Increase the counter by one if $|t| > t_{0.025}(n - 2)$. Also, add
$\vec{b}$ to the two-dimensional vector.
\end{enumerate}

\item After the do loop, divide the counter by the number of replications to
calculate the frequency of rejecting the null. Also, divide the
two-dimensional vector that has accumulated $\vec{b}$ by the number of
replications. It should equal $\E(\vec{b} \mid \vec{x})$ if the number of
replications is infinite.
\end{enumerate}

Note that in this first simulation, $\vec{x}$ is \textit{fixed} throughout
the do loop for $\vec{y}$. The second simulation calculates the
\textit{un}conditional distribution of the $t$-ratio. It should consist of
the following steps.

\begin{enumerate}[label=(\arabic*)]
\item Set the counter to zero.

\item Start a do loop of a large number of replications. In each
replication, do the following.

\begin{enumerate}[label=(\roman*)]
\item (Generate $\vec{x}$)\quad Draw a vector $\bm{\eta}$ of $n$ i.i.d.\
random variables from $N(0, 1)$ and $x_0$ from
$N(c/(1 - \phi),\; 1/(1 - \phi^2))$, and calculate $\vec{x}$
by~\eqref{eq:ps1.mc.stars}.

\item (Generate $\vec{y}$)\quad Draw a vector $\bm{\varepsilon}$ of $n$
i.i.d.\ random variables from $N(0, 1)$, and calculate
$\vec{y} = (y_1, \ldots, y_n)'$ by~($*$).

\item From a sample $(\vec{y}, \vec{x})$ thus generated, calculate the
$t$-value for $\mathrm{H}_0\colon \beta_2 = 0.5$ from the sample
$(\vec{y}, \vec{x})$.

\item Increase the counter by one if $|t| > t_{0.025}(n - 2)$.
\end{enumerate}

\item After the do loop, divide the counter by the number of replications.
\end{enumerate}

For the two simulations, verify that, for a sufficiently large number of
replications,

\begin{enumerate}[label=\arabic*.]
\item the mean of $\vec{b}$ from the first simulation is arbitrarily close
to the true value $(1, 0.5)$;

\item the frequency of rejecting the true hypothesis $\mathrm{H}_0$ (the
type~I error) is arbitrarily close to 5 percent in either simulation.
\end{enumerate}

\item In those two simulations, is the (nonconstant) regressor strictly
exogenous? Is the error conditionally homoskedastic?
\end{enumerate}

\end{problemset}

%% ====================================================================

\begin{answers}

\subsection*{Analytical Exercises}

\answer{1}
$(\vec{y} - \vec{X}\widetilde{\bm{\beta}})'
 (\vec{y} - \vec{X}\widetilde{\bm{\beta}})$

$= [(\vec{y} - \vec{X}\vec{b}) + \vec{X}(\vec{b} - \widetilde{\bm{\beta}})]'
   [(\vec{y} - \vec{X}\vec{b}) + \vec{X}(\vec{b} - \widetilde{\bm{\beta}})]$

\hspace{2em}(by the add-and-subtract strategy)

$= [(\vec{y} - \vec{X}\vec{b})' + (\vec{b} - \widetilde{\bm{\beta}})'\vec{X}']
   [(\vec{y} - \vec{X}\vec{b}) + \vec{X}(\vec{b} - \widetilde{\bm{\beta}})]$

$= (\vec{y} - \vec{X}\vec{b})'(\vec{y} - \vec{X}\vec{b})
 + (\vec{b} - \widetilde{\bm{\beta}})'\vec{X}'(\vec{y} - \vec{X}\vec{b})$

$\quad + (\vec{y} - \vec{X}\vec{b})'\vec{X}(\vec{b} - \widetilde{\bm{\beta}})
 + (\vec{b} - \widetilde{\bm{\beta}})'\vec{X}'\vec{X}
   (\vec{b} - \widetilde{\bm{\beta}})$

$= (\vec{y} - \vec{X}\vec{b})'(\vec{y} - \vec{X}\vec{b})
 + 2(\vec{b} - \widetilde{\bm{\beta}})'\vec{X}'(\vec{y} - \vec{X}\vec{b})
 + (\vec{b} - \widetilde{\bm{\beta}})'\vec{X}'\vec{X}
   (\vec{b} - \widetilde{\bm{\beta}})$

\hspace{2em}(since $(\vec{b} - \widetilde{\bm{\beta}})'\vec{X}'
(\vec{y} - \vec{X}\vec{b})
= (\vec{y} - \vec{X}\vec{b})'\vec{X}(\vec{b} - \widetilde{\bm{\beta}})$)

$= (\vec{y} - \vec{X}\vec{b})'(\vec{y} - \vec{X}\vec{b})
 + (\vec{b} - \widetilde{\bm{\beta}})'\vec{X}'\vec{X}
   (\vec{b} - \widetilde{\bm{\beta}})$

\hspace{2em}(since $\vec{X}'(\vec{y} - \vec{X}\vec{b}) = \vec{0}$
by the normal equations)

$\geq (\vec{y} - \vec{X}\vec{b})'(\vec{y} - \vec{X}\vec{b})$

\hspace{2em}(since $(\vec{b} - \widetilde{\bm{\beta}})'\vec{X}'\vec{X}
(\vec{b} - \widetilde{\bm{\beta}}) = \vec{z}'\vec{z}
= \sum_{i=1}^{n} z_i^2 \geq 0$ where
$\vec{z} \equiv \vec{X}(\vec{b} - \widetilde{\bm{\beta}})$).

\bigskip

\answer{7a} $\widehat{\bm{\beta}}_{\text{GLS}} - \bm{\beta} = \vec{A}\bm{\varepsilon}$
where $\vec{A} \equiv (\vec{X}'\vec{V}^{-1}\vec{X})^{-1}\vec{X}'\vec{V}^{-1}$
and $\vec{b} - \widehat{\bm{\beta}}_{\text{GLS}} = \vec{B}\bm{\varepsilon}$
where
$\vec{B} \equiv (\vec{X}'\vec{X})^{-1}\vec{X}'
- (\vec{X}'\vec{V}^{-1}\vec{X})^{-1}\vec{X}'\vec{V}^{-1}$.
So
%
\begin{align*}
\Cov(\widehat{\bm{\beta}}_{\text{GLS}} - \bm{\beta},\,
     \vec{b} - \widehat{\bm{\beta}}_{\text{GLS}})
&= \Cov(\vec{A}\bm{\varepsilon},\, \vec{B}\bm{\varepsilon}) \\
&= \vec{A}\,\Var(\bm{\varepsilon})\,\vec{B}' \\
&= \sigma^2 \vec{A}\vec{V}\vec{B}'.
\end{align*}
%
It is straightforward to show that $\vec{A}\vec{V}\vec{B}' = \vec{0}$.

\bigskip

\answer{7b} For the choice of $\vec{H}$ indicated in the hint,
%
\[
\Var(\widehat{\bm{\beta}}) - \Var(\widehat{\bm{\beta}}_{\text{GLS}})
= -\vec{C}\vec{V}_q^{-1}\vec{C}'.
\]
%
If $\vec{C} \neq \vec{0}$, then there exists a nonzero vector $\vec{z}$
such that $\vec{C}'\vec{z} \equiv \vec{v} \neq \vec{0}$. For such $\vec{z}$,
%
\[
\vec{z}'\bigl[\Var(\widehat{\bm{\beta}})
- \Var(\widehat{\bm{\beta}}_{\text{GLS}})\bigr]\vec{z}
= -\vec{v}'\vec{V}_q^{-1}\vec{v} < 0
\quad \text{(since $\vec{V}_q$ is positive definite)},
\]
%
which is a contradiction because $\widehat{\bm{\beta}}_{\text{GLS}}$ is
efficient.

\bigskip

\subsection*{Empirical Exercises}

\answer{a} Nerlove's description in Appendix~B.4 leads one to believe that
he did not include the depreciation rate $\delta$ in his construction of the
price of capital.

\medskip

\answer{b} Your estimates should agree with~(1.7.7).

\medskip

\answer{c} Our estimates differ from Nerlove's slightly. This would happen
even if the data used by Nerlove were the same as those provided to you,
because computers in his age were much less precise and had more frequent
rounding errors.

\medskip

\answer{d} How well can you replicate Nerlove's reported results? Fairly
well. The point estimates of returns to scale in each of the five subsamples
are 2.5, 1.5, 1.1, 1.1, and .96. As the level of output increases, the
returns to scale decline.

\medskip

\answer{e} Model~2 can be written as $\vec{y} = \vec{X}\bm{\beta} + \bm{\varepsilon}$,
where $\vec{y}$, $\vec{X}$, and $\bm{\varepsilon}$ are as
in~\eqref{eq:ps1.emp.star}. So (setting $j = 2$),
%
\[
\vec{X}'\vec{X}
= \begin{bmatrix}
\vec{X}^{(1)\prime}\vec{X}^{(1)} & \vec{0} \\
\vec{0} & \vec{X}^{(2)\prime}\vec{X}^{(2)}
\end{bmatrix},
\]
%
which means
%
\[
(\vec{X}'\vec{X})^{-1}
= \begin{bmatrix}
\bigl(\vec{X}^{(1)\prime}\vec{X}^{(1)}\bigr)^{-1} & \vec{0} \\
\vec{0} & \bigl(\vec{X}^{(2)\prime}\vec{X}^{(2)}\bigr)^{-1}
\end{bmatrix}.
\]
%
And
%
\[
\vec{X}'\vec{y}
= \begin{bmatrix}
\vec{X}^{(1)\prime}\vec{y}^{(1)} \\
\vec{X}^{(2)\prime}\vec{y}^{(2)}
\end{bmatrix}.
\]
%
Therefore,
%
\[
(\vec{X}'\vec{X})^{-1}\vec{X}'\vec{y}
= \begin{bmatrix}
\bigl(\vec{X}^{(1)\prime}\vec{X}^{(1)}\bigr)^{-1}
\vec{X}^{(1)\prime}\vec{y}^{(1)} \\[4pt]
\bigl(\vec{X}^{(2)\prime}\vec{X}^{(2)}\bigr)^{-1}
\vec{X}^{(2)\prime}\vec{y}^{(2)}
\end{bmatrix}.
\]
%
Thus, the OLS estimate of the coefficient vector for Model~2 is the same as
that for Model~1. Since the estimate of the coefficient vector is the same,
the sum of squared residuals, too, is the same.

\medskip

\answer{f} The number of restrictions is 16. $K = \text{\#coefficients in
Model~2} = 20$. So the two degrees of freedom should be $(16, 125)$.
$SSR_U = 12.262$ and $SSR_R = 21.640$. $F$-ratio $= 5.97$ with a $p$-value
of 0.0000. So this can be rejected at any reasonable significance level.

\medskip

\answer{g} $SSR_U = 12.262$ and $SSR_R = 12.577$. So $F = .40$ with 8 and
125 degrees of freedom. Its $p$-value is 0.92. So the restrictions can be
accepted at any reasonable significance level. Nerlove's $F$-ratio (see
p.~183, 8th line from bottom) is 1.576.

\medskip

\answer{h} The plot still shows that the conditional second moment is
somewhat larger for smaller firms, but now there is no evidence for possible
nonlinearities.

\end{answers}
